\section{Quick-reference grammar}

This is a practical description of the accepted surface language rather than a
separate formal language specification:
\begin{lstlisting}[language={}]
protocol     := [params-line] logical-line*
params-line  := "PARAMS:" (name ":" type)*
logical-line := command | comment | blank

command      := opcode argument*
opcode       := supported-command
             | primitive-gate
             | primitive-gate ("dg" | "inv")
             | ("dg" | "inv") primitive-gate
             | "PHASE=" angle
             | "PHASE" ("dg" | "inv") "=" angle
             | ("dg" | "inv") "PHASE=" angle

gate-row     := opcode "-I" input+ "-O" output+ ["-$R"]
                ["-IF" name "=" ("0" | "1")]
measure-row  := "MEASURE" ["-I" input]
set-row      := "SET" target value
process-row  := "MAIN-PROCESS" name
child-row    := ("RUNCHILD" | "CALL") name ["-DEPTH" integer]
                ["-I" argument*] ["-O" argument*]
rec-row      := ("REC" | "TREC") ["MAXDEPTH" integer]
exit-row     := "EXIT" "WHEN" ("DEPTH" | "LEVEL" | "ROOTDEPTH")
                ("==" | "!=" | "<" | "<=" | ">" | ">=") integer
recur-row    := "RECUR" "-I" argument+ ["-O" argument*]
if-block     := "IF" name "=" ("0" | "1")
                logical-line*
                ["ELSE" logical-line*]
                "ENDIF"
returns-row  := "RETURNVALS" token*
accept-row   := "ACCEPTVALS" token*

token        := ["$"] name [":" suffix]
constant     := ("0p" | "1p" | "sp") ["_dim" integer]
angle        := rational
             | rational "d"
             | [sign] [integer ["*"]] "pi" ["/" integer]
\end{lstlisting}

\textbf{Syntax detail:} \code{PHASE= pi/2} is joined into one opcode/value pair,
but \code{PHASE = pi/2} is not. Operations and flags are case-insensitive;
names, types, and positional values are otherwise retained.

\textbf{Block and recursion detail:} \code{-DEPTH} is required when the child
declares \code{REC} or \code{TREC}. \code{RECUR}, \code{REC}, \code{TREC},
and \code{EXIT WHEN} are only meaningful inside a child process. An
\code{if-block} lowers each inner gate to a \code{-IF} condition; see
\hyperlink{if-else-blocks}{Structured classical branches}.
